\documentclass[11pt,a4paper]{article}
\usepackage[margin=25mm]{geometry}
\usepackage[T1]{fontenc}
\usepackage{lmodern,microtype}
\usepackage{amsmath,amssymb}
\usepackage{amsthm}
\usepackage{needspace}
\theoremstyle{definition}
\newtheorem{definition}{Definition}[section]
\usepackage[hidelinks]{hyperref}
\usepackage[backend=biber,style=numeric]{biblatex}
\addbibresource{macros/AdaptiveSteane.bib}

\usepackage{macros/tikzfig}
\input{macros/floquet.tikzstyles}
\input{macros/floquet.tikzdefs}

\title{Adaptive Steane error correction without postselection}
\author{}
\date{}

\begin{document}
\maketitle

\noindent
Steane error correction measures syndromes by applying transversal CNOTs
between the data and encoded ancillas, then measuring the ancillas. In a
postselected preparation scheme, an ancilla is discarded if any verification flag is
nonzero, and preparation is repeated until all flags are zero. Here we use
every prepared ancilla and adapt the subsequent extraction circuits and
decoding to the observed flags and syndromes. Our goal is a QEC cycle
satisfying Gottesman's fault-tolerance conditions
(Appendix~\ref{sec:gottesman-ft}) without postselection.

We consider Steane error correction for an $[[n,k,d]]$ CSS code, with
\[
  t=\left\lfloor\frac{d-1}{2}\right\rfloor.
\]
\section{One Steane extraction}
\label{sec:extraction}

Let $H_Z$ and $H_X$ be the binary check matrices of the Z and X
stabilizers. Write $X(w)=\prod_i X_i^{w_i}$ and
$Z(w)=\prod_i Z_i^{w_i}$, with all vector sums taken modulo two.
The incoming data error and the ancilla error immediately before the
transversal CNOTs are
\[
    E_D=X(u)Z(v),\qquad E_A=X(x)Z(z),
    \qquad u,v,x,z\in\mathbb F_2^n.
\]
The following equations take the transversal CNOTs and final measurements
to be ideal.

Z-SE uses the encoded all-plus state $\ket{\overline{+^k}}$ as the CNOT
target and measures the ancilla in the Z basis. X-SE uses the encoded all-zero state
$\ket{\overline{0^k}}$ as the control and measures the ancilla in the X basis.

\begin{figure}[htbp]
    \centering
    \begin{minipage}{0.48\textwidth}
        \centering
        \tikzfig{figures/steane-z-extraction}
    \end{minipage}\hfill
    \begin{minipage}{0.48\textwidth}
        \centering
        \tikzfig{figures/steane-x-extraction}
    \end{minipage}
    \caption{Alternative single-basis extractions. Doubled wires represent
    $n$-qubit blocks, and each CNOT is transversal. Preparation returns
    flags $f$ and an ancilla with error $E_A$. The final measurement is
    processed into the syndrome $\sigma_Z$ or $\sigma_X$.}
    \label{fig:steane-extractions}
\end{figure}

Before any correction, the measured syndrome and output data error are
\begin{equation}
    \begin{array}{c|c|c}
        \text{Extraction} & \text{Syndrome} & E_D' \\
        \hline
        Z\text{-SE} & \sigma_Z=H_Z(u+x) & X(u)Z(v+z) \\
        X\text{-SE} & \sigma_X=H_X(v+z) & X(u+x)Z(v).
    \end{array}
    \label{eq:extraction}
\end{equation}
In Z-SE, the ancilla X error changes the measurement outcome, while its Z error
propagates to the data. In X-SE, these roles are exchanged. We retain
both the preparation flags and the measured syndrome for the adaptive
step.

\paragraph{Transversal CNOT and measurement faults.}
A bit flip in a Z-basis measurement is equivalent to an ancilla X error
just before measurement; a bit flip in an X-basis measurement is
equivalent to an ancilla Z error.
These errors commute backwards through the respective CNOTs, so they can
be placed at the preparation output without changing its flags.

Model a faulty physical CNOT by a Pauli error after the ideal gate. On
its data--ancilla pair, write
\[
    P=X_D^a Z_D^b X_A^c Z_A^e,
    \qquad a,b,c,e\in\mathbb F_2.
\]
The following replacements give the same measurement outcomes and
conditional data state. The effective ancilla error is placed before
the CNOT and the residual data error after it:
\begin{equation}
    \begin{array}{c|c|c}
        \text{Extraction} & \text{Effective ancilla error}
                         & \text{Residual data error} \\
        \hline
        Z\text{-SE} & X_A^c Z_A^b & X_D^a \\
        X\text{-SE} & X_A^a Z_A^e & Z_D^b.
    \end{array}
    \label{eq:coupling-faults}
\end{equation}
This follows from CNOT Pauli propagation, omitting a final ancilla Z in
Z-SE or X in X-SE, which does not affect the measurement. The two entries
in a row arise from one physical fault, even when both are nonidentity;
their joint effects must be retained when checking ECRP and ECCP
(Appendix~\ref{sec:gottesman-ft}).

The residual data X from Z-SE passes through a subsequent ideal X-SE
without changing its syndrome and remains a single data error. The dual
statement holds for a residual data Z from X-SE followed by ideal Z-SE.

\section{Adaptive Steane error correction}
\label{sec:adaptive}

\subsection{Preparation guarantees}
\label{sec:preparation-guarantees}

For each encoded ancilla state $\ket A$, we use a finite library of
preparation circuits with different circuit distances $\delta$; the
data-code distance remains $d$.
A Pauli fault set $F$ produces flags $f(F)$ and ancilla error $E_F$.
Flags have fault-free value zero.
In the chosen physical fault model, assume
\begin{equation}
    \begin{gathered}
        f(F)=f(F'),\qquad \sigma(E_F)=\sigma(E_{F'}),\qquad
        |F|+|F'|<\delta \\
        \Longrightarrow\qquad
        E_F\ket A\propto E_{F'}\ket A.
    \end{gathered}
    \label{eq:preparation-distance}
\end{equation}
Here $\sigma(X(x)Z(z))=(H_Zx,H_Xz)$ is the ideal code syndrome, and
$b_\delta=\lfloor(\delta-1)/2\rfloor$. Thus flags and a perfect
ancilla syndrome determine a correction through $b_\delta$ preparation
faults, including nonzero flags.

We also require that zero flags imply no increase in error weight:
\begin{equation}
    \begin{gathered}
        f(F)=0,\quad |F|\leq b_\delta
        \quad\Longrightarrow\quad
        E_F\ket A\propto\widetilde E_F\ket A,\\
        \operatorname{wt}(\widetilde E_F)\leq |F|
        \quad\text{for some Pauli }\widetilde E_F.
    \end{gathered}
    \label{eq:quiet-preparation}
\end{equation}
The weight is the ordinary Pauli weight used in the appendix. Zero flags
therefore allow faults, but exclude error amplification within the
preparation's protection order.
The logical-completion proof of Eq.~\eqref{eq:quiet-preparation} requires
that every violating zero-flag pattern $F$ admit $q_{\rm add}$ additional
single-qubit data faults producing an undetectable logical change of
$\ket A$, with $|F|+q_{\rm add}<\delta$. If
Eq.~\eqref{eq:preparation-distance} includes these added faults, comparison
with $F'=\varnothing$ gives a contradiction. This completion property
must be established for the chosen family; otherwise verify the weight
bound directly.

\Needspace{13\baselineskip}
\subsection{Adaptive decoding}
\label{sec:adaptive-decoding}

The measured Steane syndrome contains both data and ancilla errors
(Eq.~\eqref{eq:extraction}). After $m$ extractions, use the complete record
$R=((B_j,\delta_j,f_j,\sigma_j))_{j=1}^{m}$, where $B_j\in\{Z,X\}$.
Write $r$ for the data-error weight at the cycle input and $s$ for the
accumulated internal fault count. For logical correctness (ECCP), retain every
possibility consistent with $R$ and $r+s\leq t$. The output-error
condition (ECRP) must also be checked on arbitrary-input possibilities
with $s\leq t$ (Appendix~\ref{sec:gottesman-ft}). Track $r$, $s$, and the
current data error in each case.
Data errors introduced by faulty transversal CNOTs are already counted
in $s$: each faulty physical location counts once, with its joint X and
Z effects retained.
Final Pauli corrections are tracked classically in a Pauli frame, and
the EC conditions are evaluated in that frame.

\Needspace{10\baselineskip}
Start with Z-SE using a distance-$d$ preparation. After each extraction:
\begin{enumerate}
    \item Update the possible errors and fault counts using the new
    flags and syndrome. If no ECRP possibility remains, finish with the
    identity frame update.
    \item If one Pauli correction satisfies ECCP and ECRP for every
    remaining possibility in their respective ranges, update the Pauli
    frame and finish.
    \item Otherwise choose another extraction basis and preparation
    circuit, including new faults when evaluating its possible outcomes.
\end{enumerate}

Let $L(R)$ be the least internal fault count consistent with $R$ over the
full ECRP range. Then
\begin{equation}
    b(R)=t-L(R)
    \label{eq:remaining-faults}
\end{equation}
bounds further faults. An available preparation with $b_\delta\geq b(R)$,
equivalently $\delta\geq2b(R)+1$, gives sufficient protection against new
faults. The full-record test must also account for the remaining data errors:
syndromes can distinguish them even when $L(R)$ is unchanged.

\paragraph{Quiet rounds.}
For a continuation policy using Z-SE followed by X-SE in each round,
call a round quiet when neither preparation raises a flag. Check that
every quiet-round record admits a common correction under the stopping
test above, retaining the flags and syndromes from earlier rounds.
When $b(R)=0$, no further faults are allowed in the bounded-fault model,
so distance-one preparations suffice for the remaining extractions.

\section{Correctness and resource cost}
\label{sec:correctness-cost}

\subsection{Correctness}

Consider a Z-SE/X-SE round policy whose quiet-round stopping test is
verified in the chosen Pauli fault model. When $s\leq t$, every update
retains the actual input error
and fault pattern. The terminal frame update therefore satisfies ECRP,
and ECCP when $r+s\leq t$.

A continuing round has a fresh nonzero preparation flag and certifies a
new fault. If $J$ rounds are begun, $J-1\leq s$, hence $J\leq s+1\leq t+1$.
Outside this fault range, $t+1$ nonquiet rounds leave no bounded-fault
possibility, triggering the identity frame update.

\Needspace{16\baselineskip}
\subsection{Resource cost}

Let $N$ count all individual extractions, including flagged ones. For
the round policy,
\begin{equation}
    N\leq 2J\leq 2(s+1)\leq 2(t+1)
    \qquad(s\leq t).
    \label{eq:extraction-cost}
\end{equation}
The empty-set rule gives $N\leq2(t+1)$ for the other records too.
Let $c(B,\delta)$ be an additive cost of a complete extraction and $c_{\max}$ its
maximum over the available circuit library. A terminal record $R$ has cost
\begin{equation}
    C(R)=\sum_{j=1}^{N}c(B_j,\delta_j)
    \leq 2(t+1)c_{\max},
    \label{eq:circuit-cost}
\end{equation}
For CNOT cost, $c$ includes preparation and transversal CNOTs. Given noise
and input-error models, $\mathbb E[C]=\sum_R p(R)C(R)$ over terminal records.

Using the full record cannot increase the optimum cost with the same
library and objective: flags-only adaptive choices remain available.
Any strict saving must be verified for the chosen policy.

\clearpage
\appendix

\section{Gottesman's Definition of Fault-tolerance\texorpdfstring{~\cite{gottesmanIntroductionQuantumError2009,gottesmanSurvivingQuantumComputer2026}}{}}
\label{sec:gottesman-ft}

\begin{definition}
    \label{def:r-filter}
    The \emph{strict $r$-filter} for the QECC $Q$ is the projector on the subspace spanned by states of the form $E \ket{\psi}$, where $\ket{\psi} \in Q$ and $E$ is an error of weight at most $r$.
    Pictorially, I will draw an $r$-filter as:
    \[
        \tikzfig{figures/Gottesman/r-filter.tikz}
    \]
    The doubled horizontal lines in this picture represent a block of the QECC $Q$ rather than single physical qubits.
\end{definition}
\begin{definition}
    The \emph{strict ideal decoder} for the QECC $Q$ is the quantum channel that takes a state (possibly with errors) encoded in $Q$, corrects the errors, and then decodes the logical state, discarding all ancilla qubits.
    I will draw an ideal decoder as follows:
    \[
        \tikzfig{figures/Gottesman/ideal-decoder.tikz}
    \]
    The doubled horizontal line on the left of the picture represents a block of $Q$.
    The single horizontal line on the right of the picture is the logical qubit after decoding.
\end{definition}

\subsection{Fault-tolerant Gate Properties}

\begin{definition}
    The picture to represent a noisy gate gadget with exactly $s$ faults is
    \[
        \tikzfig{figures/Gottesman/noisy-1-gate.tikz}
        \qquad
        \text{or}
        \qquad
        \tikzfig{figures/Gottesman/noisy-2-gate.tikz}
    \]
    for one- and two-qubit gates, respectively.
    Each doubled horizontal line represents a single block of the QECC\@.

    The picture to represent an ideal gate performed on unencoded qubits without noise is
    \[
        \tikzfig{figures/Gottesman/ideal-1-gate.tikz}
        \qquad
        \text{or}
        \qquad
        \tikzfig{figures/Gottesman/ideal-2-gate.tikz}
    \]
    for one- and two-qubit gates, respectively. Each single horizontal line represents a single unencoded qubit.
\end{definition}

\begin{definition}[FT Gate Error Propagation Property]
    Suppose we have a gate gadget associated with a QECC with distance $2t + 1$ (i.e., it corrects $t$-errors).
    If it is a single-qubit gate gadget, it satisfies the \emph{FT Gate Error Propagation Property}, abbreviated \emph{GPP}, if, whenever $r + s \leq t$,
    \[
        \tikzfig{figures/Gottesman/GPP1.tikz}
    \]
    A two-qubit gate gadget satisfies the GPP if, whenever $r_1 + r_2 + s \leq t$,
    \[
        \tikzfig{figures/Gottesman/GPP2.tikz}
    \]
    For both equations, the faulty locations on both sides of the equation have the same fault path and error maps substituted in.
\end{definition}
\begin{definition}[FT Gate Correctness Property]
    Suppose we have a gate gadget associated with a QECC with distance $2t + 1$.
    If it is a single-qubit gate gadget, it satisfies the FT Gate Correctness Property, abbreviated GCP, if, whenever $r + s \leq t$,
    \[
        \tikzfig{figures/Gottesman/GCP1.tikz}
    \]
    A two-qubit gate gadget satisfies the GCP if, whenever $r_1 + r_2 + s \leq t$,
    \[
        \tikzfig{figures/Gottesman/GCP2.tikz}
    \]
    These equations must hold for any choice of CP maps substituted in for the faulty locations on the left-hand side.
\end{definition}

\subsection{Fault-tolerant State Preparation}

\begin{definition}
    The picture for a single-qubit state preparation gadget with exactly $s$ faults is
    \[
        \tikzfig{figures/Gottesman/noisy-state-prep.tikz}
    \]
    The picture for a single-qubit measurement gadget with exactly $s$ faults is
    \[
        \tikzfig{figures/Gottesman/noisy-measurement.tikz}
    \]
    In both cases, the doubled horizontal lines represent encoded blocks of the QECC\@.
    The measurement gadget outputs classical information, generally a single bit.

    The pictures for ideal preparation of a single-qubit state and ideal single-qubit measurement are
    \[
        \tikzfig{figures/Gottesman/ideal-state-prep.tikz}
        \qquad
        \text{and}
        \qquad
        \tikzfig{figures/Gottesman/ideal-measurement.tikz}
    \]
    The horizontal lines represent single unencoded qubits, and the ideal measurement outputs a classical bit.
\end{definition}

\begin{definition}[FT Preparation Error Propagation Property]
    Suppose we have a preparation gadget associated with a QECC with distance $2t + 1$.
    The gadget satisfies the \emph{FT Preparation Error Propagation Property}, abbreviated \emph{PPP}, if, whenever $s \leq t$,
    \[
        \tikzfig{figures/Gottesman/PPP.tikz}
    \]
\end{definition}

\begin{definition}[Preparation Correctness Property]
    Suppose we have a preparation gadget associated with a QECC with distance $2t + 1$.
    The gadget satisfies the \emph{FT Preparation Correctness Property}, abbreviated \emph{PCP}, if, whenever $s \leq t$,
    \[
        \tikzfig{figures/Gottesman/PCP.tikz}
    \]
\end{definition}

\Needspace{9\baselineskip}
\begin{definition}[FT Measurement Correctness Property]
    Suppose we have a preparation gadget associated with a QECC with distance $2t + 1$.
    The gadget satisfies the \emph{FT Measurement Correctness Property}, abbreviated \emph{MCP}, if, whenever $r + s \leq t$,
    \[
        \tikzfig{figures/Gottesman/MCP.tikz}
    \]
\end{definition}

\subsection{Fault-tolerant Quantum Error Correction}

\begin{definition}
    A fault-tolerant error correction gadget (which I will usually abbreviate EC in the future) is represented by the following picture:
    \[
        \tikzfig{figures/Gottesman/QEC-Box.tikz}
    \]
    The horizontal lines represent blocks of the QECC.
\end{definition}

\begin{definition}[FT Error Correction Recovery Property]
    Suppose we have an error correction gadget associated with a QECC with distance $2t + 1$.
    The gadget satisfies the \emph{FT Error Correction Recovery Property}, abbreviated \emph{ECRP}, if, whenever $s \leq t$,
    \[
        \tikzfig{figures/Gottesman/ECRP.tikz}
    \]
\end{definition}
\begin{definition}[FT Error Correction Correctness Property]
    Suppose we have an error correction gadget associated with a QECC with distance $2t + 1$.
    The gadget satisfies the \emph{FT Error Correction Correctness Property}, abbreviated \emph{ECCP}, if, whenever $r + s \leq t$,
    \[
        \tikzfig{figures/Gottesman/ECCP.tikz}
    \]
\end{definition}

\printbibliography

\end{document}
